% ===========================================================================
%  Programmation Frontend — Lab 3 : useEffect et Cycle de Vie
%  ESP/UCAD — 2025–2026
%  Dr. El Hadji Bassirou TOURÉ
%  Compilé avec : xelatex (2–3 passes)
% ===========================================================================
\documentclass[11pt,a4paper]{article}

\usepackage{fontspec}
\setmainfont{Latin Modern Roman}
\setsansfont{Latin Modern Sans}
\setmonofont{Latin Modern Mono}
\usepackage{polyglossia}
\setdefaultlanguage{french}

\usepackage[margin=2.5cm]{geometry}
\usepackage{amsmath,amssymb}
\usepackage{graphicx}
\usepackage{xcolor}
\usepackage{booktabs}
\usepackage{hyperref}
\usepackage{tcolorbox}
\tcbuselibrary{theorems,breakable,skins}
\usepackage{enumitem}
\usepackage{fancyhdr}
\usepackage{array}
\usepackage{pifont}

\newcommand{\cmark}{\ding{51}}
\newcommand{\xmark}{\ding{55}}

\hypersetup{colorlinks=true,linkcolor=blue!60!black,urlcolor=blue!60!black,citecolor=blue!60!black}

\newtcolorbox{objectifbox}[1][]{colback=blue!5,colframe=blue!50,fonttitle=\bfseries,title=#1,breakable}
\newtcolorbox{infobox}[1][]{colback=green!5,colframe=green!50!black,fonttitle=\bfseries,title=#1,breakable}
\newtcolorbox{attentionbox}[1][]{colback=orange!5,colframe=orange!70,fonttitle=\bfseries,title=#1,breakable}
\newtcolorbox{retenirbox}[1][]{colback=yellow!10,colframe=yellow!50!black,fonttitle=\bfseries,title=#1,breakable}
\newtcolorbox{prereqbox}[1][]{colback=cyan!5,colframe=cyan!50!black,fonttitle=\bfseries,title=#1,breakable}
\newtcolorbox{projetbox}[1][]{colback=teal!5,colframe=teal!50!black,fonttitle=\bfseries,title=#1,breakable}

\usepackage{listings}
\lstdefinelanguage{JavaScript}{keywords={break,case,catch,continue,debugger,default,delete,do,else,finally,for,function,if,in,instanceof,new,return,switch,this,throw,try,typeof,var,void,while,with,const,let,class,export,extends,import,super,yield,async,await,of,from,true,false,null,undefined},morecomment=[l]{//},morecomment=[s]{/*}{*/},morestring=[b]',morestring=[b]",morestring=[b]`,sensitive=true}
\lstdefinelanguage{CSS}{keywords={color,background,border,padding,margin,font,display,flex,align,justify,width,height,gap,cursor,none,solid,center,column,row,wrap,pointer,radius,size,weight,family,items,content,direction,transform,transition,hover,box,shadow,max,min,text,decoration,opacity},morecomment=[s]{/*}{*/},morestring=[b]',morestring=[b]",sensitive=true}
\lstset{basicstyle=\ttfamily\small,backgroundcolor=\color{gray!8},frame=single,framerule=0.5pt,rulecolor=\color{gray!40},breaklines=true,breakatwhitespace=true,showstringspaces=false,tabsize=2,xleftmargin=4pt,xrightmargin=4pt,aboveskip=8pt,belowskip=8pt,language=JavaScript,keywordstyle=\color{blue!70!black}\bfseries,commentstyle=\color{green!50!black}\itshape,stringstyle=\color{red!60!black},literate={é}{{\'e}}1 {è}{{\`e}}1 {ê}{{\^e}}1 {à}{{\`a}}1 {ù}{{\`u}}1 {ô}{{\^o}}1 {î}{{\^i}}1 {ç}{{\c{c}}}1 {É}{{\'E}}1 {→}{{$\rightarrow$}}1}
\lstdefinestyle{terminal}{language={},keywordstyle={},commentstyle={},stringstyle={},basicstyle=\ttfamily\small,backgroundcolor=\color{black!5}}
\lstdefinestyle{css}{language=CSS,keywordstyle=\color{blue!70!black}\bfseries,commentstyle=\color{green!50!black}\itshape,stringstyle=\color{red!60!black}}

\pagestyle{fancy}\fancyhf{}
\fancyhead[L]{\small Prog. Frontend --- Lab 3 : useEffect et Cycle de Vie}
\fancyhead[R]{\small ESP/UCAD}
\fancyfoot[C]{\thepage}
\fancyfoot[L]{\small Dr. El Hadji Bassirou TOURÉ}
\fancyfoot[R]{\small 2025--2026}
\renewcommand{\headrulewidth}{0.4pt}
\renewcommand{\footrulewidth}{0.4pt}

\title{\vspace{-1cm}\textbf{\LARGE Programmation Frontend}\\[0.8cm]\textbf{\huge Lab 3 --- useEffect et Cycle de Vie}\\[0.3cm]\textbf{\Large Chargement automatique, loading et page de détail}\\[0.5cm]\large ESP/UCAD --- 2025--2026}
\author{Dr. El Hadji Bassirou TOURÉ\\Département Génie Informatique\\Université Cheikh Anta Diop de Dakar}
\date{}

\begin{document}
\maketitle
\thispagestyle{empty}
\vspace{0.5cm}

\begin{prereqbox}[Informations pratiques]
\begin{center}
\begin{tabular}{lp{10cm}}
\toprule
\textbf{Durée estimée} & 2h (une séance) \\
\textbf{Prérequis} & Avoir terminé le Lab 2bis (tri, favoris, immutabilité). Comprendre \texttt{useState} et le re-rendu. \\
\textbf{Fichiers de départ} & Le projet \texttt{annuaire-react/} du Lab 2bis \\
\textbf{Livrables} & Annuaire avec chargement automatique, gestion loading/erreur, et vue de détail \\
\textbf{Validation} & Les données se chargent au démarrage. Un clic sur une carte affiche le détail. \\
\bottomrule
\end{tabular}
\end{center}
\end{prereqbox}

\begin{objectifbox}[Objectifs du lab]
À l'issue de ce lab, chaque étudiant sera capable de :
\begin{enumerate}[nosep]
\item Expliquer ce qu'est un \textbf{effet de bord} (\emph{side effect}) dans React
\item Utiliser \texttt{useEffect} pour exécuter du code \textbf{au montage} d'un composant
\item Gérer les trois états d'un chargement : \textbf{loading}, \textbf{erreur} et \textbf{données}
\item Utiliser \texttt{useEffect} avec une \textbf{dépendance} pour réagir à un changement
\item Naviguer entre une vue \textbf{liste} et une vue \textbf{détail} avec un état
\end{enumerate}
\end{objectifbox}

\tableofcontents
\newpage

% ===========================================================================
\section{Point de départ --- Rappel du Lab 2bis}
% ===========================================================================

Au Lab 2bis, vous avez compris \texttt{useState} en profondeur : le re-rendu, l'immutabilité, la différence entre état et props. L'annuaire a maintenant un tri A-Z/Z-A et un système de favoris. Votre \texttt{App.jsx} utilise 5 états :

\begin{center}
\begin{tabular}{lp{8cm}}
\toprule
\textbf{État} & \textbf{Rôle} \\
\midrule
\texttt{users} & Le tableau des utilisateurs (récupérés par fetch) \\
\texttt{chargement} & Booléen : le fetch est-il en cours ? \\
\texttt{recherche} & Le texte tapé dans la barre de recherche \\
\texttt{triAscendant} & Booléen : tri A$\to$Z ou Z$\to$A ? \\
\texttt{favoris} & Tableau d'identifiants des utilisateurs favoris \\
\bottomrule
\end{tabular}
\end{center}

\begin{objectifbox}[Étape 0 --- Vérification]
Ouvrez le projet \texttt{annuaire-react/} dans VS Code. Lancez le serveur :
\begin{lstlisting}[style=terminal]
cd annuaire-react
npm run dev
\end{lstlisting}
Ouvrez \texttt{http://localhost:5173}. Cliquez sur «~Charger~». Vérifiez que les 10 utilisateurs apparaissent, que le filtre, le tri et les favoris fonctionnent.
\end{objectifbox}

\subsection{Le problème}

Ouvrez la page dans un \textbf{nouvel onglet}. Que voyez-vous ? Un écran vide avec un bouton. Il faut cliquer sur «~Charger~» pour voir les utilisateurs.

Sur n'importe quel site réel, les données sont déjà là quand la page s'affiche. On veut la même chose : les utilisateurs doivent se charger \textbf{automatiquement} quand la page s'ouvre.

% ===========================================================================
\section{Partie 1 --- Les effets de bord}
% ===========================================================================

\begin{objectifbox}[Objectif --- 5 minutes]
Comprendre pourquoi React a besoin d'un outil séparé pour certaines actions.
\end{objectifbox}

Au Lab 2bis, vous avez vu que la fonction d'un composant a un seul rôle : \textbf{calculer ce qui doit s'afficher} à partir de l'état et des props. React l'exécute à chaque re-rendu.

Mais parfois, un composant doit aussi faire des choses \textbf{en dehors} du rendu :
\begin{itemize}[nosep]
\item Récupérer des données depuis une API (\texttt{fetch})
\item Démarrer un timer (\texttt{setInterval})
\item Modifier le titre de la page (\texttt{document.title})
\end{itemize}

Ces actions s'appellent des \textbf{effets de bord} (\emph{side effects}). On ne peut pas les mettre directement dans le corps du composant, car elles se déclencheraient à \textbf{chaque rendu} (rappelez-vous le \texttt{console.log} du Lab 2bis qui s'affichait à chaque clic).

React fournit \texttt{useEffect} pour dire : «~après le rendu, exécute ceci.~»

\begin{infobox}[Analogie : le post-it sur le frigo]
Vous vous réveillez le matin (le composant se monte). Vous voyez un post-it sur le frigo : «~Acheter du pain~». Vous faites d'abord votre routine (React fait le rendu), puis vous exécutez le post-it (l'effet).

\texttt{useEffect} est le post-it de React : «~Après le rendu, fais ceci.~»
\end{infobox}

% ===========================================================================
\section{Partie 2 --- useEffect : charger au montage}
% ===========================================================================

\begin{objectifbox}[Objectif --- 20 minutes]
Utiliser \texttt{useEffect} pour charger les utilisateurs automatiquement et supprimer le bouton.
\end{objectifbox}

\subsection{La syntaxe}

\begin{lstlisting}
import { useState, useEffect } from "react";

useEffect(() => {
  // Code à exécuter après le rendu
  console.log("Le composant a été rendu !");
}, []);
\end{lstlisting}

\texttt{useEffect} prend deux arguments : une \textbf{fonction} (le code à exécuter) et un \textbf{tableau de dépendances} (quand l'exécuter).

\begin{center}
\begin{tabular}{lp{8cm}}
\toprule
\textbf{Dépendances} & \textbf{Quand l'effet se déclenche} \\
\midrule
\texttt{[]} (tableau vide) & Une seule fois, au \textbf{montage} (quand le composant apparaît) \\
\texttt{[x]} & Au montage + à chaque fois que \texttt{x} change \\
\texttt{[x, y]} & Au montage + à chaque fois que \texttt{x} ou \texttt{y} change \\
Pas de tableau & À \textbf{chaque} rendu (rarement utile) \\
\bottomrule
\end{tabular}
\end{center}

Pour charger les données au démarrage, on veut que l'effet se déclenche \textbf{une seule fois}. On utilise donc \texttt{[]}.

\subsection{Exercice 1 --- Supprimer le bouton}

\begin{objectifbox}[À vous de jouer]
Remplacez le bouton «~Charger~» par un \texttt{useEffect} qui charge les données automatiquement.
\end{objectifbox}

Quatre changements à faire dans \texttt{src/App.jsx} :

\begin{enumerate}[nosep]
\item Ajoutez \texttt{useEffect} à l'import : \texttt{import \{ useState, useEffect \} from "react";}
\item Déplacez la logique de \texttt{chargerUtilisateurs} dans un \texttt{useEffect}
\item Initialisez \texttt{chargement} à \texttt{true} (on charge dès le départ)
\item Supprimez le \texttt{<button>} du JSX
\end{enumerate}

Voici le haut du fichier (le reste ne change pas) :

\begin{lstlisting}
import { useState, useEffect } from "react";
import UserCard from "./components/UserCard";
import SearchBar from "./components/SearchBar";
import styles from "./App.module.css";

const App = () => {
  const [users, setUsers] = useState([]);
  const [chargement, setChargement] = useState(true);
  const [recherche, setRecherche] = useState("");
  const [triAscendant, setTriAscendant] = useState(true);
  const [favoris, setFavoris] = useState([]);

  useEffect(() => {
    const chargerUtilisateurs = async () => {
      try {
        const response = await fetch(
          "https://jsonplaceholder.typicode.com/users"
        );
        const data = await response.json();
        setUsers(data);
      } catch (error) {
        console.error("Erreur :", error);
      }
      setChargement(false);
    };

    chargerUtilisateurs();
  }, []);

  // toggleFavori, usersFiltres... inchangés
\end{lstlisting}

Dans le \texttt{return}, supprimez le bloc du bouton «~Charger~». Le reste (SearchBar, barre de tri, favoris, cartes) reste identique.

Sauvegardez. Ouvrez la page. Les utilisateurs apparaissent \textbf{automatiquement}.

\begin{attentionbox}[Pourquoi la fonction async est dans le useEffect]
\texttt{useEffect} ne peut pas être \texttt{async} directement. On définit donc une fonction async \textbf{à l'intérieur} et on l'appelle immédiatement :
\begin{lstlisting}
useEffect(() => {
  const fetchData = async () => {
    // ... await fetch(...)
  };
  fetchData();
}, []);
\end{lstlisting}
\end{attentionbox}

\begin{retenirbox}[Avant / après : le chargement]
\begin{center}
\begin{tabular}{lp{5cm}p{5cm}}
\toprule
& \textbf{Lab 2bis} & \textbf{Lab 3} \\
\midrule
\textbf{Déclencheur} & Bouton «~Charger~» & Automatique (\texttt{useEffect}) \\
\textbf{Code} & \texttt{onClick=\{fn\}} & \texttt{useEffect(() => \{...\}, [])} \\
\textbf{Premier affichage} & Écran vide & Données visibles \\
\bottomrule
\end{tabular}
\end{center}
\end{retenirbox}

% ===========================================================================
\section{Partie 3 --- Gérer le loading et les erreurs}
% ===========================================================================

\begin{objectifbox}[Objectif --- 15 minutes]
Afficher un message de chargement et gérer les erreurs proprement.
\end{objectifbox}

\subsection{Le pattern loading / erreur / données}

Quand on récupère des données depuis une API, trois situations sont possibles :

\begin{center}
\begin{tabular}{lp{4cm}p{5cm}}
\toprule
\textbf{Situation} & \textbf{Ce qui se passe} & \textbf{Ce qu'on affiche} \\
\midrule
\textbf{Loading} & La requête est en cours & «~Chargement en cours...~» \\
\textbf{Erreur} & La requête a échoué & «~Impossible de charger~» \\
\textbf{Données} & La requête a réussi & La liste des utilisateurs \\
\bottomrule
\end{tabular}
\end{center}

Ce pattern est \textbf{universel} : vous le retrouverez dans chaque application qui consomme une API.

\subsection{Exercice 2 --- Ajouter un état d'erreur}

\begin{objectifbox}[À vous de jouer]
Ajoutez un état \texttt{erreur} et affichez un message conditionnel selon la situation.
\end{objectifbox}

Ajoutez l'état et modifiez le \texttt{useEffect} :

\begin{lstlisting}
const [erreur, setErreur] = useState(null);

useEffect(() => {
  const chargerUtilisateurs = async () => {
    try {
      const response = await fetch(
        "https://jsonplaceholder.typicode.com/users"
      );
      if (!response.ok) {
        throw new Error("Erreur serveur");
      }
      const data = await response.json();
      setUsers(data);
    } catch (error) {
      setErreur(error.message);
    }
    setChargement(false);
  };

  chargerUtilisateurs();
}, []);
\end{lstlisting}

La ligne \texttt{if (!response.ok)} vérifie le code HTTP. Si le serveur répond avec une erreur (404, 500...), on lance une exception.

Ajoutez deux \emph{early returns} \textbf{avant} le \texttt{return} principal :

\begin{lstlisting}
if (chargement) {
  return (
    <div className={styles.container}>
      <p>Chargement en cours...</p>
    </div>
  );
}

if (erreur) {
  return (
    <div className={styles.container}>
      <p className={styles.erreur}>
        Erreur : {erreur}
      </p>
    </div>
  );
}

// return principal (la liste) inchangé
\end{lstlisting}

Ajoutez la classe dans \texttt{src/App.module.css} :

\begin{lstlisting}[style=css]
.erreur {
  color: #c0392b;
  font-weight: bold;
  padding: 1rem;
  background: #fdecea;
  border-radius: 6px;
}
\end{lstlisting}

\begin{retenirbox}[Le pattern early return]
\begin{enumerate}[nosep]
\item \textbf{Si loading} $\to$ afficher «~Chargement...~» et \texttt{return}.
\item \textbf{Si erreur} $\to$ afficher le message d'erreur et \texttt{return}.
\item \textbf{Sinon} $\to$ afficher les données.
\end{enumerate}
On traite d'abord les cas particuliers, puis le cas général.
\end{retenirbox}

\subsection{Tester l'erreur}

Pour vérifier, changez temporairement l'URL en \texttt{/FAUX}. Le message d'erreur rouge apparaît. Remettez \texttt{/users} après le test.

% ===========================================================================
\section{Partie 4 --- Vue de détail avec dépendance}
% ===========================================================================

\begin{objectifbox}[Objectif --- 20 minutes]
Cliquer sur une carte pour afficher les informations complètes d'un utilisateur.
\end{objectifbox}

\subsection{Le principe}

On veut deux vues : la \textbf{liste} (ce qu'on a) et le \textbf{détail} (informations complètes). Un état \texttt{selectedUserId} contrôle la vue active : \texttt{null} = liste, un identifiant = détail.

\begin{infobox}[Pourquoi pas un vrai routeur ?]
Le routing (changement d'URL, bouton retour du navigateur) est un sujet à part entière. On l'apprendra au \textbf{Lab 4}. Aujourd'hui, on se concentre sur \textbf{une seule chose} : \texttt{useEffect}. La navigation par état est une solution temporaire.
\end{infobox}

\subsection{Créer \texttt{UserDetail.jsx}}

Créez \texttt{src/components/UserDetail.jsx} :

\begin{lstlisting}
import { useState, useEffect } from "react";
import styles from "./UserDetail.module.css";

const UserDetail = ({ userId, onRetour }) => {
  const [user, setUser] = useState(null);
  const [chargement, setChargement] = useState(true);
  const [erreur, setErreur] = useState(null);

  useEffect(() => {
    const chargerDetail = async () => {
      setChargement(true);
      setErreur(null);
      try {
        const response = await fetch(
          `https://jsonplaceholder.typicode.com/users/${userId}`
        );
        if (!response.ok) {
          throw new Error("Utilisateur introuvable");
        }
        const data = await response.json();
        setUser(data);
      } catch (error) {
        setErreur(error.message);
      }
      setChargement(false);
    };

    chargerDetail();
  }, [userId]);

  if (chargement) return <p>Chargement du profil...</p>;

  if (erreur) {
    return (
      <div>
        <p className={styles.erreur}>
          Erreur : {erreur}
        </p>
        <button className={styles.retour}
                onClick={onRetour}>
          Retour à la liste
        </button>
      </div>
    );
  }

  return (
    <div className={styles.detail}>
      <button className={styles.retour}
              onClick={onRetour}>
        Retour à la liste
      </button>
      <h2 className={styles.nom}>{user.name}</h2>
      <table className={styles.tableau}>
        <tbody>
          <tr>
            <td className={styles.label}>Email</td>
            <td>{user.email}</td>
          </tr>
          <tr>
            <td className={styles.label}>Téléphone</td>
            <td>{user.phone}</td>
          </tr>
          <tr>
            <td className={styles.label}>Site web</td>
            <td>{user.website}</td>
          </tr>
          <tr>
            <td className={styles.label}>Entreprise</td>
            <td>{user.company.name}</td>
          </tr>
          <tr>
            <td className={styles.label}>Ville</td>
            <td>{user.address.city}</td>
          </tr>
        </tbody>
      </table>
    </div>
  );
};

export default UserDetail;
\end{lstlisting}

Deux choses à observer :
\begin{enumerate}[nosep]
\item Le \texttt{useEffect} utilise \texttt{[userId]} comme dépendance. Si \texttt{userId} change, l'effet se re-déclenche et charge les nouvelles données.
\item Le composant applique le même pattern loading/erreur/données que \texttt{App}.
\end{enumerate}

\subsection{Créer \texttt{UserDetail.module.css}}

Créez \texttt{src/components/UserDetail.module.css} :

\begin{lstlisting}[style=css]
.detail { padding: 1rem 0; }

.retour {
  padding: 0.4rem 1rem;
  font-size: 0.9rem;
  cursor: pointer;
  border: 1px solid #ccc;
  border-radius: 6px;
  background: #f9f9f9;
  margin-bottom: 1rem;
}
.retour:hover { background: #eee; }

.nom {
  margin: 0.5rem 0 1rem 0;
  color: #333;
}

.tableau { width: 100%; border-collapse: collapse; }
.tableau td {
  padding: 0.5rem 0.75rem;
  border-bottom: 1px solid #eee;
}
.label {
  font-weight: bold;
  color: #555;
  width: 120px;
}
.erreur { color: #c0392b; font-weight: bold; }
\end{lstlisting}

\subsection{Rendre les cartes cliquables}

Modifiez \texttt{src/components/UserCard.jsx} pour ajouter une prop \texttt{onClick} :

\begin{lstlisting}
import styles from "./UserCard.module.css";

const UserCard = ({
  user, estFavori, onToggleFavori, onClick
}) => {
  return (
    <div className={styles.card}
         onClick={() => onClick(user.id)}>
      <div className={styles.header}>
        <h3 className={styles.nom}>{user.name}</h3>
        <button className={styles.etoile}
                onClick={(e) => {
                  e.stopPropagation();
                  onToggleFavori(user.id);
                }}>
          {estFavori ? "\u2605" : "\u2606"}
        </button>
      </div>
      <p className={styles.info}>
        Email : {user.email}
      </p>
      <p className={styles.info}>
        Entreprise : {user.company.name}
      </p>
    </div>
  );
};

export default UserCard;
\end{lstlisting}

\begin{attentionbox}[\texttt{e.stopPropagation()} sur le bouton étoile]
Sans \texttt{e.stopPropagation()}, cliquer sur l'étoile déclencherait \textbf{aussi} le \texttt{onClick} de la carte (l'étoile est à l'intérieur). \texttt{stopPropagation} empêche l'événement de remonter. Le clic sur l'étoile reste un toggle favori, le clic sur le reste de la carte ouvre le détail.
\end{attentionbox}

Ajoutez \texttt{cursor: pointer;} dans \texttt{.card} de \texttt{UserCard.module.css} si ce n'est pas déjà fait.

\subsection{Connecter le tout dans \texttt{App.jsx}}

Ajoutez l'état \texttt{selectedUserId}, l'import de \texttt{UserDetail}, et un early return pour la vue détail. Voici le fichier complet :

\begin{lstlisting}
import { useState, useEffect } from "react";
import UserCard from "./components/UserCard";
import UserDetail from "./components/UserDetail";
import SearchBar from "./components/SearchBar";
import styles from "./App.module.css";

const App = () => {
  const [users, setUsers] = useState([]);
  const [chargement, setChargement] = useState(true);
  const [erreur, setErreur] = useState(null);
  const [recherche, setRecherche] = useState("");
  const [triAscendant, setTriAscendant] = useState(true);
  const [favoris, setFavoris] = useState([]);
  const [selectedUserId, setSelectedUserId] =
    useState(null);

  useEffect(() => {
    const chargerUtilisateurs = async () => {
      try {
        const response = await fetch(
          "https://jsonplaceholder.typicode.com/users"
        );
        if (!response.ok) {
          throw new Error("Erreur serveur");
        }
        const data = await response.json();
        setUsers(data);
      } catch (error) {
        setErreur(error.message);
      }
      setChargement(false);
    };

    chargerUtilisateurs();
  }, []);

  const toggleFavori = (userId) => {
    if (favoris.includes(userId)) {
      setFavoris(favoris.filter(id => id !== userId));
    } else {
      setFavoris([...favoris, userId]);
    }
  };

  if (chargement) {
    return (
      <div className={styles.container}>
        <p>Chargement en cours...</p>
      </div>
    );
  }

  if (erreur) {
    return (
      <div className={styles.container}>
        <p className={styles.erreur}>
          Erreur : {erreur}
        </p>
      </div>
    );
  }

  if (selectedUserId) {
    return (
      <div className={styles.container}>
        <h1 className={styles.titre}>
          Annuaire des utilisateurs
        </h1>
        <UserDetail
          userId={selectedUserId}
          onRetour={() => setSelectedUserId(null)}
        />
      </div>
    );
  }

  const usersFiltres = users
    .filter(user =>
      user.name.toLowerCase().includes(
        recherche.toLowerCase()
      )
    )
    .sort((a, b) => {
      if (triAscendant) {
        return a.name.localeCompare(b.name);
      }
      return b.name.localeCompare(a.name);
    });

  return (
    <div className={styles.container}>
      <h1 className={styles.titre}>
        Annuaire des utilisateurs
      </h1>

      <SearchBar
        recherche={recherche}
        onRecherche={setRecherche}
      />

      <div className={styles.barre}>
        <p className={styles.compteur}>
          {usersFiltres.length} utilisateur(s)
        </p>
        <button className={styles.boutonTri}
                onClick={() =>
                  setTriAscendant(!triAscendant)
                }>
          Tri : {triAscendant ? "A → Z" : "Z → A"}
        </button>
      </div>

      {favoris.length > 0 && (
        <p className={styles.favoris}>
          {favoris.length} favori(s)
        </p>
      )}

      {usersFiltres.map(user => (
        <UserCard
          key={user.id}
          user={user}
          estFavori={favoris.includes(user.id)}
          onToggleFavori={toggleFavori}
          onClick={setSelectedUserId}
        />
      ))}
    </div>
  );
};

export default App;
\end{lstlisting}

\subsection{Tester}

Sauvegardez tous les fichiers. Vérifiez :
\begin{enumerate}[nosep]
\item Les utilisateurs se chargent \textbf{automatiquement}.
\item Cliquez sur une carte $\to$ le détail s'affiche.
\item Cliquez sur «~Retour à la liste~» $\to$ la liste réapparaît.
\item Les étoiles de favoris fonctionnent toujours (l'étoile ne déclenche pas la vue détail).
\item Le filtre et le tri fonctionnent toujours.
\end{enumerate}

% ===========================================================================
\section{Bilan --- Structure du projet}
% ===========================================================================

\begin{lstlisting}[style=terminal]
annuaire-react/
  src/
    App.jsx
    App.module.css
    components/
      UserCard.jsx
      UserCard.module.css
      UserDetail.jsx        <-- nouveau
      UserDetail.module.css  <-- nouveau
      SearchBar.jsx
      SearchBar.module.css
    main.jsx
  index.html
  package.json
\end{lstlisting}

\begin{retenirbox}[Deux useEffect --- deux rôles]

\begin{center}
\begin{tabular}{lp{3.5cm}p{5cm}}
\toprule
\textbf{Composant} & \textbf{Dépendances} & \textbf{Rôle} \\
\midrule
\texttt{App} & \texttt{[]} & Charger la liste une seule fois au démarrage \\
\texttt{UserDetail} & \texttt{[userId]} & Charger le détail quand l'id change \\
\bottomrule
\end{tabular}
\end{center}
\end{retenirbox}

\begin{retenirbox}[Avant / après : Lab 2bis vs Lab 3]
\begin{center}
\begin{tabular}{lp{4.5cm}p{4.5cm}}
\toprule
& \textbf{Lab 2bis} & \textbf{Lab 3} \\
\midrule
\textbf{Chargement} & Bouton «~Charger~» & Automatique (\texttt{useEffect}) \\
\textbf{Loading} & Texte dans le bouton & Écran dédié (early return) \\
\textbf{Erreurs} & \texttt{console.error} & Message visible à l'écran \\
\textbf{Détail} & Aucun & Vue complète au clic \\
\textbf{Navigation} & Aucune & Liste $\leftrightarrow$ Détail par état \\
\bottomrule
\end{tabular}
\end{center}
\end{retenirbox}

% ===========================================================================
\section{Aide-mémoire}
% ===========================================================================

\begin{center}
\begin{tabular}{lp{8cm}}
\toprule
\textbf{Concept} & \textbf{Syntaxe / Exemple} \\
\midrule
Importer useEffect & \texttt{import \{ useState, useEffect \} from "react";} \\
Au montage seulement & \texttt{useEffect(() => \{ ... \}, []);} \\
Quand \texttt{x} change & \texttt{useEffect(() => \{ ... \}, [x]);} \\
À chaque rendu & \texttt{useEffect(() => \{ ... \});} (rare) \\
\midrule
Async dans useEffect & Définir une fonction async à l'intérieur, puis l'appeler \\
Vérifier la réponse & \texttt{if (!response.ok) throw new Error(...);} \\
\midrule
Early return (loading) & \texttt{if (chargement) return <p>...</p>;} \\
Early return (erreur) & \texttt{if (erreur) return <p>...</p>;} \\
\midrule
Navigation par état & \texttt{const [id, setId] = useState(null);} \\
Vue conditionnelle & \texttt{if (id) return <Detail />;} \\
Bloquer la propagation & \texttt{e.stopPropagation()} \\
\bottomrule
\end{tabular}
\end{center}

\newpage

% ===========================================================================
\section{Résumé et conclusion}
% ===========================================================================

Dans ce lab, vous avez :
\begin{enumerate}[nosep]
\item Compris les \textbf{effets de bord} : des actions en dehors du rendu (fetch, timers).
\item Supprimé le bouton et chargé les données \textbf{automatiquement} avec \texttt{useEffect(\ldots, [])}.
\item Mis en place le pattern \textbf{loading / erreur / données} avec des \emph{early returns}.
\item Créé un composant \texttt{UserDetail} avec \texttt{useEffect} et une \textbf{dépendance} \texttt{[userId]}.
\item Navigué entre liste et détail grâce à un état \texttt{selectedUserId}.
\end{enumerate}

\begin{center}
\begin{tabular}{lcp{7cm}}
\toprule
\textbf{Étape} & \textbf{Statut} & \textbf{Ce que vous savez faire} \\
\midrule
\textbf{DOM et événements} & \cmark{} & Sélecteurs, événements, modification du DOM \\
\textbf{Fonctions JS} & \cmark{} & Déclaration, expression, fléchée \\
\textbf{Asynchrone} & \cmark{} & \texttt{fetch}, \texttt{async/await} \\
\textbf{React --- les bases} & \cmark{} & Composant, JSX, Vite, \texttt{.map()} \\
\textbf{Composants et props} & \cmark{} & Découpage, props, CSS Modules \\
\textbf{État React} & \cmark{} & \texttt{useState}, re-rendu, immutabilité \\
\textbf{useEffect} & \cmark{} & Montage, dépendances, loading/erreur \\
\textbf{Formulaires et routing} & $\to$ Lab 4 & Inputs contrôlés, React Router \\
\bottomrule
\end{tabular}
\end{center}

\begin{projetbox}[Et ensuite ?]
L'annuaire charge automatiquement, gère les erreurs et permet de voir le détail d'un utilisateur. Mais la navigation est gérée par un état : l'URL ne change pas et le bouton retour du navigateur ne fonctionne pas.

Dans le \textbf{Lab 4}, vous installerez \textbf{React Router} pour créer une vraie application multi-pages. Chaque vue aura sa propre URL (\texttt{/users}, \texttt{/users/3}), et vous ajouterez un \textbf{formulaire contrôlé} pour saisir de nouveaux utilisateurs.
\end{projetbox}

% ===========================================================================
\section*{Exercice bonus --- Pour aller plus loin}
% ===========================================================================

\begin{infobox}[Bonus : un hook personnalisé useFetch]
Extrayez la logique fetch/loading/erreur dans un hook réutilisable. Créez \texttt{src/hooks/useFetch.js} :

\begin{lstlisting}
import { useState, useEffect } from "react";

const useFetch = (url) => {
  const [data, setData] = useState(null);
  const [loading, setLoading] = useState(true);
  const [error, setError] = useState(null);

  useEffect(() => {
    const fetchData = async () => {
      setLoading(true);
      setError(null);
      try {
        const response = await fetch(url);
        if (!response.ok) throw new Error("Erreur");
        const result = await response.json();
        setData(result);
      } catch (err) {
        setError(err.message);
      }
      setLoading(false);
    };

    fetchData();
  }, [url]);

  return { data, loading, error };
};

export default useFetch;
\end{lstlisting}

Utilisation dans \texttt{UserDetail.jsx} :

\begin{lstlisting}
import useFetch from "../hooks/useFetch";

const UserDetail = ({ userId, onRetour }) => {
  const { data: user, loading, error } = useFetch(
    `https://jsonplaceholder.typicode.com/users/${userId}`
  );

  if (loading) return <p>Chargement...</p>;
  if (error) return <p>Erreur : {error}</p>;

  return (
    <div>
      <h2>{user.name}</h2>
      {/* ... le reste du JSX */}
    </div>
  );
};
\end{lstlisting}

Le hook encapsule le pattern loading/erreur/données. Il est réutilisable dans n'importe quel composant qui charge des données.
\end{infobox}

\vspace{1cm}

\begin{center}
\fbox{
\begin{minipage}{0.85\textwidth}
\centering
\vspace{0.5cm}
\textbf{\Large Lab 3 --- useEffect et Cycle de Vie --- Terminé \cmark}\\[0.3cm]
\normalsize
Vous savez charger des données automatiquement,\\
gérer les états de chargement et afficher une vue de détail.\\[0.2cm]
Prochaine étape : React Router et formulaires contrôlés.\\[0.3cm]
\textit{useEffect = le post-it de React.}\\[0.5cm]
\end{minipage}
}
\end{center}

\end{document}
